% Algebra 2 Lesson 4-2 fast-track packet: Graphing Rational Functions (enVision A2, Example 2 and Example 3 only)
% Re-cut 2026-10-31 toward the Topic 4 test (a2_T4_assessment_TE.tex): keep Examples 2, 3 (was 3, 4).
% Keep/cut is a proposal; the PLC has not met on this lesson.
% Shared preamble for the "fast" (hacked-down) Algebra 2 lesson packets, Oct 2026.
% Each packet: % Shared preamble for the "fast" (hacked-down) Algebra 2 lesson packets, Oct 2026.
% Each packet: % Shared preamble for the "fast" (hacked-down) Algebra 2 lesson packets, Oct 2026.
% Each packet: \input{fast-packet-template} then \begin{document} ... \end{document}.
\documentclass[12pt]{article}
\usepackage[margin=0.7in]{geometry}
\usepackage{amsmath,amssymb}
\usepackage{booktabs,tabularx,array}
\usepackage{enumitem}
\usepackage[T1]{fontenc}
\usepackage[utf8]{inputenc}
\usepackage{lmodern}
\usepackage{tikz}
\usepackage{pgfplots}
\pgfplotsset{compat=1.18}
\usepackage{xcolor}
\setlist{nosep,leftmargin=1.4em}
\setlength{\parindent}{0pt}
\setlength{\parskip}{6pt}
\renewcommand{\arraystretch}{1.35}
\newcolumntype{Y}{>{\raggedright\arraybackslash}X}
\pagestyle{empty}

% blanks and workspace
\newcommand{\blank}[1][2.6cm]{\rule[-2pt]{#1}{0.5pt}}
\newcommand{\lines}[1][1]{\par\vspace{2pt}\foreach \i in {1,...,#1}{\rule{0pt}{22pt}\rule{\linewidth}{0.4pt}\par}}
\newcommand{\wspart}[2]{\vspace{4pt}{\large\bfseries #1: #2}\par\vspace{-6pt}\rule{\textwidth}{0.4pt}\par}

% header: \packethead{Lesson 1-3}{Piecewise-Defined Functions}
\newcommand{\packethead}[2]{%
{\LARGE\bfseries #2}\hfill Name \blank[5cm]\\[2pt]
Algebra 2 \quad #1 (fast track) \hfill Date \blank[2.5cm] \quad Period \blank[1.2cm]\par\vspace{2pt}}

% boxed goals: \goals{learning goal}{language goal}
\newcommand{\goals}[2]{%
\fbox{\parbox{\dimexpr\textwidth-2\fboxsep-2\fboxrule}{%
\textbf{Learning goal:} #1\\[2pt]
\textbf{Language goal:} #2}}\par}

% quiz target box: \quiztarget{text}
\newcommand{\quiztarget}[1]{%
\vspace{4pt}\colorbox{gray!12}{\parbox{\dimexpr\textwidth-2\fboxsep}{\textbf{On the quiz:} #1}}\par}

% IXL table row helper: code & skill & which quiz question it feeds
\newcommand{\ixlrow}[3]{\texttt{#1} & #2 & #3 \\}

% one grid style for every graph; override xmin/xmax/... per plot
\pgfplotsset{sheetgrid/.style={
  axis lines=middle, grid=both, grid style={gray!45},
  xmin=-10.5, xmax=10.5, ymin=-10.5, ymax=10.5,
  xtick={-10,-8,...,10}, ytick={-10,-8,...,10},
  x=0.28cm, y=0.28cm,
  ticklabel style={font=\scriptsize, fill=white, inner sep=1pt},
  xlabel={$x$}, ylabel={$y$},
  axis line style={thick}, clip=true}}
% unit grid for small windows
\pgfplotsset{unitgrid/.style={sheetgrid,
  xmin=-5.5, xmax=5.5, ymin=-5.5, ymax=5.5,
  xtick={-5,...,5}, ytick={-5,...,5}, x=0.5cm, y=0.5cm}}
% open and closed endpoints
\tikzset{closed/.style={circle, fill, inner sep=1.6pt},
         open/.style={circle, draw, fill=white, inner sep=1.4pt}}
 then \begin{document} ... \end{document}.
\documentclass[12pt]{article}
\usepackage[margin=0.7in]{geometry}
\usepackage{amsmath,amssymb}
\usepackage{booktabs,tabularx,array}
\usepackage{enumitem}
\usepackage[T1]{fontenc}
\usepackage[utf8]{inputenc}
\usepackage{lmodern}
\usepackage{tikz}
\usepackage{pgfplots}
\pgfplotsset{compat=1.18}
\usepackage{xcolor}
\setlist{nosep,leftmargin=1.4em}
\setlength{\parindent}{0pt}
\setlength{\parskip}{6pt}
\renewcommand{\arraystretch}{1.35}
\newcolumntype{Y}{>{\raggedright\arraybackslash}X}
\pagestyle{empty}

% blanks and workspace
\newcommand{\blank}[1][2.6cm]{\rule[-2pt]{#1}{0.5pt}}
\newcommand{\lines}[1][1]{\par\vspace{2pt}\foreach \i in {1,...,#1}{\rule{0pt}{22pt}\rule{\linewidth}{0.4pt}\par}}
\newcommand{\wspart}[2]{\vspace{4pt}{\large\bfseries #1: #2}\par\vspace{-6pt}\rule{\textwidth}{0.4pt}\par}

% header: \packethead{Lesson 1-3}{Piecewise-Defined Functions}
\newcommand{\packethead}[2]{%
{\LARGE\bfseries #2}\hfill Name \blank[5cm]\\[2pt]
Algebra 2 \quad #1 (fast track) \hfill Date \blank[2.5cm] \quad Period \blank[1.2cm]\par\vspace{2pt}}

% boxed goals: \goals{learning goal}{language goal}
\newcommand{\goals}[2]{%
\fbox{\parbox{\dimexpr\textwidth-2\fboxsep-2\fboxrule}{%
\textbf{Learning goal:} #1\\[2pt]
\textbf{Language goal:} #2}}\par}

% quiz target box: \quiztarget{text}
\newcommand{\quiztarget}[1]{%
\vspace{4pt}\colorbox{gray!12}{\parbox{\dimexpr\textwidth-2\fboxsep}{\textbf{On the quiz:} #1}}\par}

% IXL table row helper: code & skill & which quiz question it feeds
\newcommand{\ixlrow}[3]{\texttt{#1} & #2 & #3 \\}

% one grid style for every graph; override xmin/xmax/... per plot
\pgfplotsset{sheetgrid/.style={
  axis lines=middle, grid=both, grid style={gray!45},
  xmin=-10.5, xmax=10.5, ymin=-10.5, ymax=10.5,
  xtick={-10,-8,...,10}, ytick={-10,-8,...,10},
  x=0.28cm, y=0.28cm,
  ticklabel style={font=\scriptsize, fill=white, inner sep=1pt},
  xlabel={$x$}, ylabel={$y$},
  axis line style={thick}, clip=true}}
% unit grid for small windows
\pgfplotsset{unitgrid/.style={sheetgrid,
  xmin=-5.5, xmax=5.5, ymin=-5.5, ymax=5.5,
  xtick={-5,...,5}, ytick={-5,...,5}, x=0.5cm, y=0.5cm}}
% open and closed endpoints
\tikzset{closed/.style={circle, fill, inner sep=1.6pt},
         open/.style={circle, draw, fill=white, inner sep=1.4pt}}
 then \begin{document} ... \end{document}.
\documentclass[12pt]{article}
\usepackage[margin=0.7in]{geometry}
\usepackage{amsmath,amssymb}
\usepackage{booktabs,tabularx,array}
\usepackage{enumitem}
\usepackage[T1]{fontenc}
\usepackage[utf8]{inputenc}
\usepackage{lmodern}
\usepackage{tikz}
\usepackage{pgfplots}
\pgfplotsset{compat=1.18}
\usepackage{xcolor}
\setlist{nosep,leftmargin=1.4em}
\setlength{\parindent}{0pt}
\setlength{\parskip}{6pt}
\renewcommand{\arraystretch}{1.35}
\newcolumntype{Y}{>{\raggedright\arraybackslash}X}
\pagestyle{empty}

% blanks and workspace
\newcommand{\blank}[1][2.6cm]{\rule[-2pt]{#1}{0.5pt}}
\newcommand{\lines}[1][1]{\par\vspace{2pt}\foreach \i in {1,...,#1}{\rule{0pt}{22pt}\rule{\linewidth}{0.4pt}\par}}
\newcommand{\wspart}[2]{\vspace{4pt}{\large\bfseries #1: #2}\par\vspace{-6pt}\rule{\textwidth}{0.4pt}\par}

% header: \packethead{Lesson 1-3}{Piecewise-Defined Functions}
\newcommand{\packethead}[2]{%
{\LARGE\bfseries #2}\hfill Name \blank[5cm]\\[2pt]
Algebra 2 \quad #1 (fast track) \hfill Date \blank[2.5cm] \quad Period \blank[1.2cm]\par\vspace{2pt}}

% boxed goals: \goals{learning goal}{language goal}
\newcommand{\goals}[2]{%
\fbox{\parbox{\dimexpr\textwidth-2\fboxsep-2\fboxrule}{%
\textbf{Learning goal:} #1\\[2pt]
\textbf{Language goal:} #2}}\par}

% quiz target box: \quiztarget{text}
\newcommand{\quiztarget}[1]{%
\vspace{4pt}\colorbox{gray!12}{\parbox{\dimexpr\textwidth-2\fboxsep}{\textbf{On the quiz:} #1}}\par}

% IXL table row helper: code & skill & which quiz question it feeds
\newcommand{\ixlrow}[3]{\texttt{#1} & #2 & #3 \\}

% one grid style for every graph; override xmin/xmax/... per plot
\pgfplotsset{sheetgrid/.style={
  axis lines=middle, grid=both, grid style={gray!45},
  xmin=-10.5, xmax=10.5, ymin=-10.5, ymax=10.5,
  xtick={-10,-8,...,10}, ytick={-10,-8,...,10},
  x=0.28cm, y=0.28cm,
  ticklabel style={font=\scriptsize, fill=white, inner sep=1pt},
  xlabel={$x$}, ylabel={$y$},
  axis line style={thick}, clip=true}}
% unit grid for small windows
\pgfplotsset{unitgrid/.style={sheetgrid,
  xmin=-5.5, xmax=5.5, ymin=-5.5, ymax=5.5,
  xtick={-5,...,5}, ytick={-5,...,5}, x=0.5cm, y=0.5cm}}
% open and closed endpoints
\tikzset{closed/.style={circle, fill, inner sep=1.6pt},
         open/.style={circle, draw, fill=white, inner sep=1.4pt}}

\begin{document}

\packethead{Lesson 4-2}{Graphing Rational Functions}
\goals{I can find the domain and the vertical and horizontal asymptotes of a rational function from its equation.}{I can explain how I found each asymptote using the words numerator, denominator, degree, asymptote, and leading coefficient.}

\quiztarget{(1) Match a graph to its equation (full Savvas quiz only). (2) Bonus, +5: find the vertical asymptotes when a factor of the denominator cancels with the numerator (Example 2). (3) Graded, 10 points: find both asymptotes of a linear-over-linear function (Examples 2 and 3). (4) Water added to a mixture (full Savvas quiz only); the Practice block is optional preparation for it. (5) Graded, 10 points: pick the function with a given horizontal asymptote by comparing degrees and leading coefficients (Example 3).}
\topictarget{4}{Question 2 asks for a domain (Example 2). Questions 3 and 12 ask for both asymptotes when a factor cancels (Examples 2, 3). Question 18 asks which functions have horizontal asymptote $y=7$ (Example 3). Questions 8 (a remainder) and 16 (asymptotes of $\frac{ax+b}{cx+d}$) come from Examples 1 and 4, which this packet skips.}

\wspart{Example 2}{Find Multiple Vertical Asymptotes of a Rational Function}
What are the vertical asymptotes for the graph of $f(x)=\dfrac{3x-2}{x^2+7x+12}$?

\begin{minipage}[t]{0.6\textwidth}
\vspace{0pt}\raggedright
Vertical \textbf{asymptotes} can occur at the $x$-values where the function is undefined. Determine where the \textbf{denominator} of the rational function is equal to 0.

\textbf{Step 1} Set the denominator equal to zero and solve.\\[2pt]
$x^2+7x+12=0$\\
$(x+3)(x+4)=0$\\
$x+3=0$ or $x+4=0$\\
$x=-3$ or $x=-4$

The possible vertical asymptotes are $x=-3$ and $x=-4$.

\textbf{Step 2} Check the \textbf{numerator} there: $3(-3)-2=-11$ and $3(-4)-2=-14$. Neither is 0, so both are asymptotes. Graph the function to check: the graph is not defined at $x=-3$ or $x=-4$. \textbf{These lines are vertical asymptotes.}

{\small Rule (Concept Summary): $x=a$ is a vertical asymptote of $\frac{P(x)}{Q(x)}$ if $Q(a)=0$ and $P(a)\neq0$.\\ Cancel all common factors first. If $x=a$ still makes the reduced denominator 0, it is a vertical asymptote; if not, the graph has a hole at $x=a$. Either way, $a$ is not in the domain.}

\textbf{Domain} (the $x$-values where $Q(x)\neq0$): all real numbers except $-3$ and $-4$.
\end{minipage}\hfill
\begin{minipage}[t]{0.38\textwidth}
\vspace{0pt}
\begin{center}
\begin{tikzpicture}
\begin{axis}[sheetgrid, xmin=-7.5, xmax=1.5, ymin=-30.5, ymax=80.5,
  xtick={-7,...,1}, xticklabels={,$-6$,,$-4$,,$-2$,,,}, ytick={-30,-20,...,80}, yticklabels={,,,,,,,,$50$,,,},
  x=0.6cm, y=0.055cm]
\node[font=\scriptsize, fill=white, inner sep=1pt, anchor=north west] at (axis cs:0,0) {$O$};
\addplot[thick, dashed, blue, <->] coordinates {(-4,-30) (-4,80)};
\addplot[thick, dashed, blue, <->] coordinates {(-3,-30) (-3,80)};
\addplot[very thick, red, <->, domain=-7.4:-4.37, samples=200] {(3*x-2)/(x^2+7*x+12)};
\addplot[very thick, red, <->, domain=-3.787:-3.176, samples=200] {(3*x-2)/(x^2+7*x+12)};
\addplot[very thick, red, <->, domain=-2.73:1.4, samples=300] {(3*x-2)/(x^2+7*x+12)};
\node[font=\scriptsize, blue, fill=white, inner sep=1pt, anchor=east] at (axis cs:-4.1,72) {$x=-4$};
\node[font=\scriptsize, blue, fill=white, inner sep=1pt, anchor=west] at (axis cs:-2.9,72) {$x=-3$};
\node[font=\scriptsize, red, fill=white, inner sep=1pt] at (axis cs:-1.2,-20) {$f$};
\end{axis}
\end{tikzpicture}
\end{center}
\end{minipage}

\textbf{Say it:} The vertical asymptotes are $x=\blank[1.2cm]$ and $x=\blank[1.2cm]$ because the \blank[2.5cm] is 0 there and the \blank[2.5cm] is not.

\newpage

\wspart{Try It 2}{Vertical asymptotes}
Find the vertical asymptotes for each function. Graph the function to check your work.

\begin{minipage}[t]{0.48\textwidth}
\textbf{a.} $g(x)=\dfrac{5x}{x^2-x-6}$

\vspace{4pt}
denominator factored: \blank[3cm]\\[6pt]
denominator is 0 at: \blank[3cm]\\[6pt]
numerator at those $x$-values: \blank[2.2cm]\\[6pt]
vertical asymptotes: \blank[3.5cm]
\end{minipage}\hfill
\begin{minipage}[t]{0.48\textwidth}
\textbf{b.} $h(x)=\dfrac{7-x}{(x-5)(x+1)(x+3)}$

\vspace{4pt}
denominator is 0 at: \blank[3cm]\\[6pt]
numerator at those $x$-values: \blank[2.2cm]\\[6pt]
vertical asymptotes: \blank[3.5cm]
\end{minipage}

\wspart{Example 3}{Find Types of Horizontal Asymptotes}
What are the horizontal asymptotes for the graph of $f(x)=\dfrac{3x-2}{x^2+7x+12}$?

A horizontal asymptote describes \textbf{end behavior}: what $y$ does as $x\to\infty$ or $x\to-\infty$. Compare the \textbf{degree} (highest power) of the \textbf{numerator} with the degree of the \textbf{denominator}. There are three cases.

\vspace{2pt}
\begin{minipage}[t]{0.32\textwidth}
\raggedright\textbf{Case 1:} degree of numerator $<$ degree of denominator. Horizontal asymptote $y=0$.

$g(x)=\dfrac{x+4}{x^2+1}$ \quad (1 $<$ 2)

\begin{tikzpicture}
\begin{axis}[sheetgrid, xmin=-60, xmax=60, ymin=-2, ymax=6,
  xtick={-48,-24,24,48}, ytick={2,4}, minor xtick={-60,-48,...,60}, minor ytick={-2,...,6},
  minor grid style={gray!45}, x=0.043cm, y=0.4cm]
\node[font=\scriptsize, fill=white, inner sep=1pt, anchor=north west] at (axis cs:0,0) {$O$};
\addplot[very thick, red, <->, domain=-57:57, samples=401] {(x+4)/(x^2+1)};
\node[font=\scriptsize, red, fill=white, inner sep=1pt] at (axis cs:30,3) {$g$};
\end{axis}
\end{tikzpicture}

The bottom grows faster, so $g(x)\to0$.
\end{minipage}\hfill
\begin{minipage}[t]{0.32\textwidth}
\raggedright\textbf{Case 2:} degree of numerator $>$ degree of denominator. \textbf{No} horizontal asymptote.

$h(x)=\dfrac{x^2+1}{x+2}$ \quad (2 $>$ 1)

\begin{tikzpicture}
\begin{axis}[sheetgrid, xmin=-25, xmax=25, ymin=-20, ymax=15,
  xtick={-20,-10,10,20}, ytick={-10,10}, yticklabels={}, minor xtick={-25,-20,...,25}, minor ytick={-20,-15,...,15},
  minor grid style={gray!45}, x=0.1cm, y=0.1cm]
\node[font=\scriptsize, fill=white, inner sep=1pt, anchor=north west] at (axis cs:0,0) {$O$};
\node[font=\scriptsize, fill=white, inner sep=1pt, anchor=east] at (axis cs:25,10) {$10$};
\node[font=\scriptsize, fill=white, inner sep=1pt, anchor=east] at (axis cs:25,-10) {$-10$};
\addplot[very thick, red, <->, domain=-17.4:-2.35, samples=150] {(x^2+1)/(x+2)};
\addplot[very thick, red, <->, domain=-1.68:16, samples=200] {(x^2+1)/(x+2)};
\node[font=\scriptsize, red, fill=white, inner sep=1pt] at (axis cs:17,6) {$h$};
\end{axis}
\end{tikzpicture}

The top grows faster, so $y\to\pm\infty$.
\end{minipage}\hfill
\begin{minipage}[t]{0.32\textwidth}
\raggedright\textbf{Case 3:} degrees equal. Horizontal asymptote $y=$ the ratio of the \textbf{leading coefficients}.

$k(x)=\dfrac{2x^2+x+1}{x^2-1}$ \quad (2 $=$ 2)

\begin{tikzpicture}
\begin{axis}[sheetgrid, xmin=-20, xmax=20, ymin=-15, ymax=15,
  xtick={-16,-8,8,16}, ytick={-12,-6,6,12}, yticklabels={}, minor xtick={-20,-16,...,20}, minor ytick={-15,-12,...,15},
  minor grid style={gray!45}, x=0.12cm, y=0.12cm, xlabel={}]
\node[anchor=north east, inner sep=1pt, yshift=-2pt] at (axis cs:20,0) {$x$};
\node[font=\scriptsize, fill=white, inner sep=1pt, anchor=south west] at (axis cs:0,0) {$O$};
\node[font=\scriptsize, fill=white, inner sep=1pt, anchor=east] at (axis cs:20,12) {$12$};
\node[font=\scriptsize, fill=white, inner sep=1pt, anchor=east] at (axis cs:20,6) {$6$};
\node[font=\scriptsize, fill=white, inner sep=1pt, anchor=east] at (axis cs:20,-6) {$-6$};
\node[font=\scriptsize, fill=white, inner sep=1pt, anchor=east] at (axis cs:20,-12) {$-12$};
\addplot[very thick, red, <->, domain=-19:-1.1005, samples=200] {(2*x^2+x+1)/(x^2-1)};
\addplot[very thick, red, <->, domain=-0.9283:0.8617, samples=100] {(2*x^2+x+1)/(x^2-1)};
\addplot[very thick, red, <->, domain=1.2116:19, samples=200] {(2*x^2+x+1)/(x^2-1)};
\node[font=\scriptsize, red, fill=white, inner sep=1pt] at (axis cs:12,8) {$k$};
\node[font=\scriptsize, red, fill=white, inner sep=0.5pt, anchor=south] at (axis cs:-11,2.6) {$y\to2$};
\end{axis}
\end{tikzpicture}

Long division: $k(x)=2+\dfrac{x+3}{x^2-1}$,\\ so $y=\frac21=2$.
\end{minipage}

\vspace{4pt}
\textbf{Back to $f$.} $f(x)=\dfrac{3x-2}{x^2+7x+12}$: the numerator has degree 1 and the denominator has degree 2. That is Case 1, so the horizontal asymptote is $\boldsymbol{y=0}$.

\textbf{Say it:} The degree of the numerator of $f$ is \blank[1cm] and the degree of the denominator is \blank[1cm], so the horizontal asymptote is \blank[2cm] because \dots\lines[1]

\newpage

\wspart{Try It 3}{Horizontal asymptotes}
What are the horizontal asymptotes of the graph of each function?

\begin{minipage}[t]{0.32\textwidth}
\textbf{a.} $g(x)=\dfrac{2x^2+x-9}{2x-8}$

\vspace{4pt}
degree on top: \blank[1cm]\\[4pt]
degree on bottom: \blank[1cm]\\[4pt]
horizontal asymptote:\\[4pt] \blank[3.5cm]
\end{minipage}\hfill
\begin{minipage}[t]{0.32\textwidth}
\textbf{b.} $h(x)=\dfrac{x^2+5x+4}{3x^2-12}$

\vspace{4pt}
degree on top: \blank[1cm]\\[4pt]
degree on bottom: \blank[1cm]\\[4pt]
horizontal asymptote:\\[4pt] \blank[3.5cm]
\end{minipage}\hfill
\begin{minipage}[t]{0.32\textwidth}
\textbf{c.} $k(x)=\dfrac{x}{(2x-1)(x+6)}$

\vspace{4pt}
degree on top: \blank[1cm]\\[4pt]
degree on bottom: \blank[1cm]\\[4pt]
horizontal asymptote:\\[4pt] \blank[3.5cm]
\end{minipage}

\textbf{Say it:} In part b the asymptote is not $y=0$ because the degrees are \blank[2cm], so I use the ratio of the \blank[3.5cm].

% R5: structural parallel of 4-2 Lesson Quiz item 4 (water added to a 12% salt mixture, f(x)=1.2/(x+10));
% no Try It rehearses that context. Same structure, different numbers and setting.
\wspart{Practice}{Adding water to a mixture}
A juice maker adds $x$ liters of water to a tank that holds 5 liters of a drink that is 20\% sugar. The fraction of the new mixture that is sugar is $f(x)=\dfrac{1}{x+5}$.

\begin{minipage}[t]{0.48\textwidth}
\textbf{a.} Why is the numerator 1?\lines[2]

\textbf{b.} How many liters of water must be added to make a 4\% sugar mixture? (4\% $=0.04$. Solve $\frac{1}{x+5}=0.04$, or graph $y=f(x)$ and $y=0.04$ and find where they meet.)\lines[3]
\end{minipage}\hfill
\begin{minipage}[t]{0.48\textwidth}
\textbf{c.} What is the horizontal asymptote of $f$? What does it mean for the drink when a lot of water is added?\lines[3]

\textbf{d.} The vertical asymptote is $x=-5$. Why does it not matter here?\lines[2]
\end{minipage}

\newpage

\wspart{IXL}{Practice for the quiz}
Go to IXL, Algebra 2. Type the code in the search box. Do the starred skills first.

\vspace{2pt}
\begin{tabularx}{\textwidth}{l Y Y}
\toprule
\textbf{Code} & \textbf{Skill} & \textbf{Feeds} \\
\midrule
\ixlrow{7JJ $\bigstar$}{Rational functions: asymptotes and excluded values}{Quiz 2, 3 and 5: both asymptotes from the equation (Examples 2 and 3)}
\ixlrow{79Q $\bigstar$}{Domain of rational functions}{Quiz 2 and 3: the denominator is 0 at a vertical asymptote and the value is left out of the domain (Example 2)}
\ixlrow{ULD}{Domain and range of reciprocal functions}{Warm-up from Lesson 4-1 (Quiz 1, full Savvas quiz only)}
\ixlrow{UB5 $\bigstar$}{Factor quadratics}{Quiz 2 (bonus): factor a denominator to find where it is 0 (Example 2, Try It 2a, Exit question)}
\bottomrule
\end{tabularx}

\vspace{4pt}
\textbf{SmartScore 70 or above on each starred ($\bigstar$) skill replaces your quiz score.}

% DOK 3: Practice & Problem Solving 10 (Error Analysis), tied to Example 3 in the TE item analysis.
\wspart{Exit question}{Describe and correct the error}
\fbox{\parbox{\dimexpr\textwidth-2\fboxsep-2\fboxrule}{%
Juanita is trying to determine the vertical and horizontal asymptotes for the graph of the function $f(x)=\dfrac{x^2+3x-4}{x^2-x-12}$. Here is her work:

\begin{center}
\begin{tabular}{l}
$f(x)=\dfrac{x^2+3x-4}{x^2-x-12}=\dfrac{(x+4)(x-1)}{(x+3)(x-4)}$\\[8pt]
vertical asymptote: $x=-3$, $x=4$\\
horizontal asymptote: $y=-4$, $y=1$ \quad\textcolor{red}{\Large$\times$}
\end{tabular}
\end{center}

Describe and correct the error Juanita made in determining the vertical and horizontal asymptotes.

\vspace{4pt}
Juanita's vertical asymptotes are \blank[2.5cm] because the \blank[2.5cm] is 0 at $x=-3$ and $x=4$ and the numerator is not.

\vspace{6pt}
Her horizontal asymptotes are wrong because she used \blank[5cm].

\vspace{6pt}
The degree of the numerator is \blank[1.2cm] the degree of the denominator, so the horizontal asymptote is $y=\blank[1.2cm]$, the ratio of the \blank[4cm].
\lines[2]
}}

\end{document}
